\honest{Failing to Learn from History}{Eliezer Yudkowsky}{2007}

When I gave my mysterious answer, I made many mistakes, but one was the most critical: I did not realize that ``solving a mystery should make it feel less confusing.'' I wanted to give the phenomenon a cause and fit it into my model of the world; why should that make it less mysterious, when mystery was its nature? \nb{The rule catches an answer that still feels mysterious. It does not catch one that feels like an explanation and is not, which is the error Yudkowsky had described the day before in ``Say Not Complexity.''}

I knew the stories of astrologers, alchemists and vitalists. But my mystery, I thought, was new and stranger, one that science had failed to explain for centuries. As if stars, matter and life had not been mysteries for thousands of years.

We learn science in school and think its subjects were never mysterious, so each generation doubts that science can explain the next great puzzle. \nb{The only case given is the author's.}

I took the lesson of history to be that the astrologers, alchemists and vitalists had a character flaw, a taste for mystery that made them give mysterious explanations of subjects that were not mysterious. Surely a truly weird phenomenon might need a weird explanation? Only later, when I began to see the ordinary structure inside my own mystery, did I see how reasonable vitalism had seemed ``at the time,'' and how surprising the universe's answer was: life is ordinary and needs no weird explanation. \nb{The history is not told in this post; the vitalist Lord Kelvin is quoted in ``Mysterious Answers to Mysterious Questions,'' five days earlier.}

We read history, but we do not live it. Had I personally postulated the old mysteries and then discovered the answers, I would have looked at my own answer and said: ``No way am I falling for that again.'' \nb{Kelvin did live through a scientific revolution in his own field, though not the dissolving of a mystery. He gave up Carnot's theory of heat for Joule's and helped formulate the laws of thermodynamics. In the passage Yudkowsky quotes, he still held the influence of life to be ``infinitely beyond the range of any scientific inquiry hitherto entered on.''}
